\section{Limits of Neumann and harmonic boundary transfer}
\label{sec:boundary195}
This section consolidates the analytic negative tests of source units
192--194. It concerns scalar finite positive-start energies and actual
rectangles. The positive partial-invariance construction above is unchanged.
None of the results below refutes the unrestricted rectangle-count question.

\subsection{Exact boundary energy}
Let a connected finite graph have positive edge conductances and energy
$E(v)=\sum_{\{i,j\}}c_{ij}(v_i-v_j)^2=v^TLv$, each edge counted once.
For a nonempty proper retained set $R$ write $B=V\setminus R$ and
$L_{RR}=L_{\mathrm{int}}+D_{\mathrm{cut}}$, where $D_{\mathrm{cut}}$ is
the diagonal of cut conductances at retained vertices.
\begin{proposition}\label{prop:boundary195-schur}
The grounded block $L_{BB}$ is positive definite, and
\begin{align*}
 \min_{v|_R=u}E(v)&=E_{\mathrm{int}}(u)+u^TQ_Ru,\\
 Q_R&=D_{\mathrm{cut}}-L_{RB}L_{BB}^{-1}L_{BR}.
\end{align*}
The form $Q_R$ is positive semidefinite and annihilates constants.
\end{proposition}
\begin{proof}
Extend a vector on $B$ by zero on $R$. If its energy vanishes,
connectedness makes it constant and hence zero. This proves positive
definiteness of $L_{BB}$. For $h=-L_{BB}^{-1}L_{BR}u$, completion of the
square gives
\[
 E(u,z)=u^T\Lambda_Ru+(z-h)^TL_{BB}(z-h),\qquad
 \Lambda_R=L_{RR}-L_{RB}L_{BB}^{-1}L_{BR}.
\]
Subtracting the fixed internal energy before minimization leaves the
nonnegative sum of cut and $B$-internal edge energies. Thus $Q_R\succeq0$.
A constant trace has a constant zero-energy extension, so $Q_R1_R=0$.
\end{proof}
This is the classical Schur or Dirichlet-to-Neumann framework
\cite{bd:kron}. The identity derives the correction but supplies no
quantitative moment bound: the inverse depends on the entire eliminated
graph. Merely naming $Q_R$ does not remove a hypothesis of a spatial theorem.

\subsection{Internal energy does not control exact extension}
On $X=(\mathbb Z/n\mathbb Z)^D$, $D\ge3$, $n\ge3$, use
\[
 E_a(v)=\sum_{x,i}a(x)|v(x+e_i)-v(x)|^2,\quad N=n^D,\quad p=D/2,
 \quad \mu=N^{-1}\sum_xa(x).
\]
Positive starts define unique unordered edges at these sizes. Internal
Neumann energy retains only edges with both endpoints in the rectangle.
For $H>1$, set $a(0)=1$ and $a(x)=H$ otherwise. Then
\[
 \mu=H-(H-1)/N,\qquad H/2\le\mu<H,\qquad 1\le a\le2\mu,
 \quad M_a=\frac{\mu^p}{N}+(1-N^{-1})(\mu/H)^p.
\]
For $R=\{0,e_1\}$ the sole internal edge has coefficient one. A mean-zero
trace $(1,-1)$ has norm squared two and internal energy four, hence sharp
local Neumann gap two.
\begin{proposition}\label{prop:boundary195-single}
For every exact extension $v$ of that trace,
$E_a(v)\ge4+4H/3$. Thus no size-independent internal-energy-only exact
extension estimate with constant $C(D,M_0)$ holds for every admissible
profile, for any fixed $M_0>1$.
\end{proposition}
\begin{proof}
The path $0\to-e_2\to e_1-e_2\to e_1$ uses three distinct strong edges,
with positive starts $-e_2$ (directions two and one) and $e_1-e_2$
(direction two). Cauchy--Schwarz gives energy at least $4H/3$ on this
path. The disjoint direct weak edge contributes four. For $H=n$,
\[
 1<M_a<1+n^{-p}\longrightarrow1.
\]
Indeed $H^p/N=n^{-p}$ bounds its first term, while its second is below
one; strict Jensen gives the lower bound. Every fixed $M_0>1$ admits
an unbounded tail, but the extension-to-internal-energy ratio is at least
$1+n/3$. This argument includes nonlinear extension selection rules.
\end{proof}
The proposition does not include $M_0=1$, when strict convexity forces
constant coefficients. Bounds allowed to depend on size, mean or a growing
target scale are different assertions.

For completeness the two-terminal conductance also satisfies
\[
 \Lambda_R=g\begin{pmatrix}1&-1\\-1&1\end{pmatrix},\qquad
 1+H/3\le g\le D(H+1),\qquad
 Q_R=(g-1)\begin{pmatrix}1&-1\\-1&1\end{pmatrix}.
\]
Symmetry and zero row sums give the displayed form of $\Lambda_R$.
The direct edge plus bypass prove the lower bound for arbitrary endpoint
values by the same telescoping calculation. For the upper bound, extend
the value at $e_1$ constantly to all vertices other than zero. Only the
$D$ outgoing weight-one edges and $D$ incoming weight-$H$ edges at zero
contribute. They are distinct for $n\ge3$, proving the upper bound.
Thus the boundary-correction to internal-energy ratio is at least $H/3$;
for $H=n$ it too has no uniform bound for any fixed $M_0>1$.

\paragraph{A contrasting uniformly coercive family.}
For the different choice $H=n^2$, $H^p=N$, so $M_a<2$ and $M_a\to2$;
the fixed admissibility bound is $M=2$, not a fixed smaller bound. Here
$\max\{1,N/(1+\mu)^p\}\le2^p$, while the global mean-zero gap is at least
four. To see the latter directly, split unit-coefficient energy as
$E_0=E_w+E_s$, the weak starts being the $D$ edges from zero. Choose one
cycle $\sigma$ through all $D$ coordinate directions and replace edge
$(0,i)$ by
\[
 0\to-e_{\sigma(i)}\to e_i-e_{\sigma(i)}\to e_i.
\]
Every replacement start is strong. For a fixed positive edge $(y,k)$,
its possible weak sources in the first, middle and last segment are
\[
 y+e_k,\qquad y+e_{\sigma(k)},\qquad
 y-e_{\sigma^{-1}(k)}+e_k.
\]
They are pairwise distinct modulo $n$ for $D\ge3,n\ge3$: the three
coordinate indices are distinct. Only source zero is weak, so an edge is
used at most once. Hence $E_w\le3E_s$, $E_0\le4E_s$ and
$E_a\ge HE_0/4$. Nonzero torus eigenvalues are at least
$4\sin^2(\pi/n)\ge16/n^2$, proving
\[
 E_a(v)\ge4\|v-\operatorname{mean}_Xv\|^2.
\]
Thus for $c=3$ one global normalized mean with penalty coefficient $K=c$
suffices, despite the displayed local gap being two. For vector fields
that cell carries one mean per component, not one scalar output in total.
The near-one family $H=n$ instead has only the comparison lower bound
$4/n$, and its target scale grows like $n^p$. These limits must not be
interchanged. No local coexact witness is asserted.

\subsection{Disjoint traces do not justify a sum of harmonic costs}
The next example refutes the aggregate estimate
\begin{equation}\label{eq:boundary195-false-sum}
 \sum_j (f|_{R_j})^TQ_{R_j}(f|_{R_j})\le C(D,M)E_a(f),
\end{equation}
where each form separately eliminates the entire complement of its own
rectangle. Joint elimination is not the same construction.

First let $T=(\mathbb Z/n\mathbb Z)^d$, $d\ge3$, with even $n$ and unit
energy $E_T$. For all $y,z\in T$,
\begin{equation}\label{eq:boundary195-point}
 |w(y)-w(z)|^2\le C_d E_T(w),\qquad C_d=d3^{d-1}/2.
\end{equation}
Here is a complete finite proof. Use the normalized Fourier coefficients
\[
 \widehat w(k)=n^{-d}\sum_y w(y)e^{-2\pi i k\cdot y/n}.
\]
Parseval gives $E_T=n^d\sum_k\lambda_k|\widehat w(k)|^2$.
Choose representatives $-n/2\le k_l<n/2$. Then
$\lambda_k=4\sum_l\sin^2(\pi k_l/n)\ge16|k|^2/n^2$.
Fourier inversion and Cauchy--Schwarz bound the point-difference constant by
\[
 n^{-d}\sum_{k\ne0}
 \frac{|e^{2\pi i k\cdot y/n}-e^{2\pi i k\cdot z/n}|^2}{\lambda_k}
 \le\frac{n^{2-d}}4\sum_{k\ne0}|k|^{-2}.
\]
The shell with $\max_l|k_l|=r$ has at most
\[
 (2r+1)^d-(2r-1)^d\le2d3^{d-1}r^{d-1}
\]
representatives. Its inverse-square contribution is at most
$2d3^{d-1}r^{d-3}$. Summing $1\le r\le n/2$ gives at most
$2d3^{d-1}n^{d-2}$, proving \eqref{eq:boundary195-point}.
The restriction $d\ge3$ is essential to this uniform estimate.

\begin{theorem}\label{thm:boundary195-aggregate}
For every fixed $D\ge4$, there are positive profiles with fixed critical
moment bound $M=3$, disjoint actual rectangles, and one global mean-zero
field for which the ratio of the left side of
\eqref{eq:boundary195-false-sum} to $E_a(f)$ tends to infinity.
\end{theorem}
\begin{proof}
Take even $n\ge6$, $d=D-1$, $p=D/2$, $m=n/2$, $H=n^{2/D}$ and
$X=(\mathbb Z/n\mathbb Z)\times T$, writing vertices $(t,y)$.
Put $a=1$ on the two planes $t=0,m$, except $a(0,0)=H$; put $a=H$
elsewhere. There are $W=2n^d-1$ weak vertices. Since $H^p=n$ and
$W/n^D<2/n\le1/3$,
\[
 \tfrac23H\le\mu<H,\quad1\le a\le2\mu,\qquad
 1\le M_a=\mu^p[W/n^D+(1-W/n^D)H^{-p}]<3.
\]
For each $y\ne0$ let $R_y=\{(0,y),(1,y)\}$. These rectangles are
disjoint and their internal coefficient is one. For arbitrary $v$,
telescope the endpoint difference through
\[
 (0,y)\to(n-1,y)\to(n-1,0)\to(0,0)\to(1,0)\to(1,y).
\]
The first, third and fourth changes use distinct longitudinal edges of
weight $H$. The two other changes are within the entirely strong
transverse layers $n-1$ and one. Their squared differences are bounded by
$C_d$ times the respective layer energy divided by $H$, by
\eqref{eq:boundary195-point}. These five energy sets are pairwise disjoint
within this one comparison: $y\ne0$ distinguishes the two incoming edges,
and layers and directions distinguish all others. Weighted Cauchy--Schwarz
with weights $1,C_d,1,1,C_d$ gives, for $K_d=3+2C_d$,
\[
 |v(0,y)-v(1,y)|^2\le K_d E_a(v)/H.
\]
The minimized two-point energy is $g_y$ times the squared endpoint
difference, so $g_y\ge H/K_d$ and its boundary correction is $g_y-1$
times that squared difference.

Let $f=1$ on layers $0,m+1,\ldots,n-1$ and $f=-1$ on layers
$1,\ldots,m$. The two sets each contain $m$ layers, so its mean is zero.
Only longitudinal starts at zero and $m$ contribute. One has weight $H$
and the remaining $2n^d-1$ have weight one, giving exactly
\[
 E_a(f)=4(2n^d-1+H),\qquad
 \sum_{y\ne0}(f|_{R_y})^TQ_{R_y}(f|_{R_y})
 \ge4(n^d-1)(H/K_d-1).
\]
Eventually $H\ge2K_d$. Since $H\le n^d$ and $n^d-1\ge n^d/2$, the
ratio is at least $H/(12K_d)$ and diverges. All constants depend only on
the fixed dimension. The moments tend to three, not a fixed smaller bound.
\end{proof}

The lower-bound certificates reuse the hole and transverse layers across
different trace problems. This does not force the actual minimizers to
use that hole rather than parallel weak cuts. Disjoint retained vertices
do not make independently optimized global energies disjoint.
The original target count remains comparable to the number of pairs:
\[
 2^{-p}n^d\le\frac{n^D}{(1+\mu)^p}\le(3/2)^p n^d.
\]
Thus the example is not a counting counterexample. There is no $D=3$
conclusion from this proof and no obstruction to a joint reduction,
partition-dependent selection, overlap compensation or a growing-factor
bound. No coexact, rank or physical continuum result is obtained.

\subsection{Attribution, assessment and unresolved mechanism}
Weighted extension theorems have their own weight and domain hypotheses;
the constants in Chua's Theorem 1.4 explicitly depend on the weights
\cite{bd:chua}. Our finite counterexamples do not contradict it.
The retained-resistance invariance under Kron reduction in
\cite{bd:kron} concerns one retained set, not a sum over different
eliminations. Fourier, path comparison and Schur reduction are classical.
No exhaustive novelty or priority claim is made for these diagnostic
families. Full source proofs and involved audits192--194 are retained at
the immutable programme commits recorded in the source manifest. Each
configured reviewer was the same historically involved Codex
\texttt{gpt-6-astra} agent with shared history, provider/account and
workspace; served backend and reasoning effort were not attested. These
were complete source audits, not blind or independent peer reviews,
scientific replay, Lean verification, PDF assessments or numerical stars.
All limitations and author responses remain in the evidence packet.

These boundary-transfer tests alone did not settle the fully variable
fixed-moment actual-rectangle question. Section~\ref{sec:balanced223}
supplies its positive finite scalar-gradient construction by a different mechanism. Cellwise internal-only extension, a near-one-moment rescue, and
independent-cell harmonic-cost additivity have explicit falsifiers.
A joint partition or direct local internal-energy argument must derive
its count without assuming any of those failed transfers. No main-theorem
hypothesis has been removed by these negative tests. No earlier PDF score
or public Working-paper status transfers to this changed manuscript.
